\section{相关工作}
\label{sec:related}

\subsection{AI辅助代码生成}

大语言模型（LLM）在代码生成任务中展现了卓越的能力~\cite{chen2021codex,li2023starcoder}。
基于这些模型构建的工具——GitHub Copilot~\cite{bird2022taking}、Anthropic的Claude Code以及Cursor——将代码生成集成到开发工作流中，使开发者能够用自然语言描述功能并直接获得可执行代码。
近期研究探索了使用这些工具进行端到端应用开发的可能性，即非程序员完全通过自然语言构建完整的应用程序~\cite{jiang2024low}~\cite{vaithilingam2022expectation}。

\subsection{代码生成的提示工程}

AI生成代码的质量在很大程度上取决于Prompt的质量。
提示工程领域的研究已经识别出几种有效策略：提供结构化规范~\cite{white2023prompt}、使用少样本示例~\cite{brown2020language}、思维链推理~\cite{wei2022chain}，以及通过对话进行迭代细化~\cite{nijkamp2022codegen}。
然而，大多数提示工程研究关注的是\textit{技术性}Prompt——由理解代码结构的开发者编写。
本文的不同之处在于聚焦于非专业用户编写的\textit{设计性}Prompt，这些用户缺乏编程和设计词汇。

\subsection{设计系统与规范}

设计系统——由清晰标准指导的可复用组件集合——已成为专业网页开发的标准实践~\cite{kholmatova2017design}。
Figma、Design Tokens和组件库等工具使团队能够在大型项目中保持视觉一致性。
在Vibe Coding的语境中，设计规范扮演着类似的角色：它充当"唯一真相来源"，约束AI的生成行为，确保不同Prompt所构建的模块之间保持连贯性。

\subsection{Vibe Coding实践}

"Vibe Coding"一词由Andrej Karpathy于2025年初推广~\cite{karpathy2025vibe}，用来描述一种工作流程：开发者"完全顺应感觉，拥抱指数增长，甚至忘记代码的存在。"
虽然早期采用者分享了关于Vibe Coding工作流的轶事经验，但对不同Vibe Coding策略——特别是针对非专业用户——的系统性研究仍然有限。
本文通过一个真实的实现案例研究，对两种不同策略进行了对比分析。
